%%%%%%%%%%%%%%%%%%%%%%%%%%%%%%%%%%%%%%%% unilateral barriers, following \cite{NSV03}
% The curves are samples of the notional profile
%   psilo(x) = 0.85 + 0.75 exp(-(x-1.2)^2/0.45) + 0.55 exp(-(x-2.9)^2/0.5)
%   u(x)     = psilo + (2.00 + 0.06 x - psilo) S(x),  S a clamped cubic smoothstep
%              which is 0 on the contact interval [0.75,1.7] and 1 beyond x=4.2
%   u_h+w    = u(x) + 0.10 sin(55 x + 40 degrees)
%   u^top    = (u_h+w) + 0.55,   u_bot = (u_h+w) - 0.6
% Sampling the formulas here, rather than plotting them with plot[domain=...],
% keeps pgfmath out of the compile; evaluating them at every sampled point was
% the entire cost of this figure.
\begin{tikzpicture}[x=2.0cm,y=1.5cm,>={Stealth[length=5pt]}]
  \newcommand{\uhw}{plot[smooth] coordinates {
    (0.000,1.281) (0.200,1.276) (0.400,1.253) (0.600,1.309) (0.800,1.475) (1.000,1.636)
    (1.200,1.698) (1.400,1.631) (1.600,1.473) (1.800,1.306) (2.000,1.229) (2.200,1.274)
    (2.400,1.406) (2.600,1.561) (2.800,1.686) (3.000,1.766) (3.200,1.823) (3.400,1.884)
    (3.600,1.963) (3.800,2.048) (4.000,2.117) (4.200,2.152) (4.400,2.166) (4.600,2.184)
    (4.800,2.205) (5.000,2.229) (5.200,2.256) (5.400,2.285) (5.600,2.315) (5.800,2.346)
    (6.000,2.377)}}
  % the star-dilated set, as a vertical strip
  \fill[Blue!7] (0.5,0.20) rectangle (1.75,3.05);
  % the band in which u is trapped
  \fill[ForestGreen!12] plot[smooth] coordinates {
    (0.000,1.831) (0.200,1.826) (0.400,1.803) (0.600,1.859) (0.800,2.025) (1.000,2.186)
    (1.200,2.248) (1.400,2.181) (1.600,2.023) (1.800,1.856) (2.000,1.779) (2.200,1.824)
    (2.400,1.956) (2.600,2.111) (2.800,2.236) (3.000,2.316) (3.200,2.373) (3.400,2.434)
    (3.600,2.513) (3.800,2.598) (4.000,2.667) (4.200,2.702) (4.400,2.716) (4.600,2.734)
    (4.800,2.755) (5.000,2.779) (5.200,2.806) (5.400,2.835) (5.600,2.865) (5.800,2.896)
    (6.000,2.927)}
    -- plot[smooth] coordinates {
    (6.000,1.777) (5.800,1.746) (5.600,1.715) (5.400,1.685) (5.200,1.656) (5.000,1.629)
    (4.800,1.605) (4.600,1.584) (4.400,1.566) (4.200,1.552) (4.000,1.517) (3.800,1.448)
    (3.600,1.363) (3.400,1.284) (3.200,1.223) (3.000,1.166) (2.800,1.086) (2.600,0.961)
    (2.400,0.806) (2.200,0.674) (2.000,0.629) (1.800,0.706) (1.600,0.873) (1.400,1.031)
    (1.200,1.098) (1.000,1.036) (0.800,0.875) (0.600,0.709) (0.400,0.653) (0.200,0.676)
    (0.000,0.681)} -- cycle;
  % mesh, strictly-signed set, and its star dilation
  \draw[gray] (0,0.30) -- (6,0.30);
  \foreach \i in {0,1,...,24} {\draw[gray] ({0.25*\i},0.22) -- ({0.25*\i},0.38);}
  \draw[Blue,line width=2pt] (0.70,0.30) -- (1.60,0.30);
  \foreach \z in {0.75,1.0,1.25,1.5} {\fill[Blue] (\z,0.30) circle (1.8pt);}
  \node[Blue,below] at (1.125,0.14) {$\{\sigma_h>0\} \ \subset\ \ombot$};
  % obstacle
  \draw[Blue,thick] plot[smooth] coordinates {
    (0.000,0.881) (0.200,0.931) (0.400,1.031) (0.600,1.187) (0.800,1.376) (1.000,1.537)
    (1.200,1.602) (1.400,1.542) (1.600,1.394) (1.800,1.236) (2.000,1.140) (2.200,1.138)
    (2.400,1.214) (2.600,1.319) (2.800,1.392) (3.000,1.390) (3.200,1.310) (3.400,1.184)
    (3.600,1.056) (3.800,0.959) (4.000,0.899) (4.200,0.869) (4.400,0.856) (4.600,0.852)
    (4.800,0.850) (5.000,0.850) (5.200,0.850) (5.400,0.850) (5.600,0.850) (5.800,0.850)
    (6.000,0.850)};
  \node[Blue,below] at (4.400,0.856) {$\psilo$};
  % corrected discrete solution, and the two barriers built from it
  \draw[gray,dashed] \uhw;
  \begin{scope}[shift={(0,0.55)}] \draw[ForestGreen,thick] \uhw; \end{scope}
  \begin{scope}[shift={(0,-0.60)}] \draw[ForestGreen,thick] \uhw; \end{scope}
  \node[ForestGreen,above,fill=white,inner sep=1pt] at (2.400,1.956) {$u^\top$};
  \node[ForestGreen,below,fill=white,inner sep=1pt] at (2.400,0.806) {$u_\bot$};
  \node[gray,below,fill=white,inner sep=1pt] at (5.100,2.242) {$u_h+w$};
  % exact solution, heavy where it contacts the obstacle
  \draw[line width=1.2pt] plot[smooth] coordinates {
    (0.000,1.216) (0.094,1.217) (0.188,1.201) (0.281,1.181) (0.375,1.166) (0.469,1.169)
    (0.562,1.195) (0.656,1.249) (0.750,1.328)};
  \draw[line width=2pt]   plot[smooth] coordinates {
    (0.750,1.328) (0.869,1.438) (0.988,1.529) (1.106,1.586) (1.225,1.601) (1.344,1.571)
    (1.462,1.502) (1.581,1.410) (1.700,1.311)};
  \draw[line width=1.2pt] plot[smooth] coordinates {
    (1.700,1.311) (1.879,1.201) (2.058,1.185) (2.237,1.264) (2.417,1.407) (2.596,1.563)
    (2.775,1.695) (2.954,1.789) (3.133,1.858) (3.312,1.922) (3.492,1.997) (3.671,2.081)
    (3.850,2.162) (4.029,2.224) (4.208,2.252) (4.388,2.263) (4.567,2.274) (4.746,2.285)
    (4.925,2.296) (5.104,2.306) (5.283,2.317) (5.462,2.328) (5.642,2.338) (5.821,2.349)
    (6.000,2.360)};
  \node[above,fill=white,inner sep=1pt] at (3.500,2.001) {$u$};
  % the two shifts, at the right margin
  \draw[dotted,gray] (6,2.377) -- (6.42,2.377);
  \draw[dotted,gray] (6,2.927) -- (6.42,2.927);
  \draw[dotted,gray] (6,1.777) -- (6.42,1.777);
  \draw[<->,thick] (6.22,2.377) -- (6.22,2.927);
  \draw[<->,thick] (6.22,2.377) -- (6.22,1.777);
  \node[right] at (6.26,2.652) {$c^\top$};
  \node[right] at (6.26,2.077) {$c_\bot$};
\end{tikzpicture}
